\chapter{Power Series} % 幂级数
\section{Power Series and Its Convergence Radius} % 幂级数和其收敛半径

\section{Expanding Functions into Power Series} % 函数的展开为幂级数
\begin{definition}{Smooth Function}\label{def:smooth_function}
    For a function \( f(x) \) defined on an interval \( I \), 
    if \( f(x) \) is continuous, then \(f(x)\) is called \(C^{0}\) function on \( I \);
    if \( f(x) \) has continuous \(n\) (\(n \geq 1\)) derivative, then \(f(x)\) is called \(C^{n}\) function on \( I \);
    if for any \(n \in \mathbb{N}\), \(f(x)\) has continuous \(n\)-th derivative, 
    then \(f(x)\) is called \(C^{\infty}\) function on \( I \), also known as \textbf{smooth function}.
\end{definition}

\begin{definition}{(Real) Analytic Function}\label{def:real_analytic_function}
    For a function \( f(x) \) defined on an interval \( I \), 
    if for any point \( x_{0} \in I \), there exists a power series expansion of \(f(x)\) at \(x_{0}\) 
    \[
    f(x) = \sum_{n=0}^{\infty} a_{n} (x - x_{0})^{n}
    \]
    that converges to \(f(x)\) in some neighborhood of \(x_{0}\), 
    then \(f(x)\) is called an \textbf{(real) analytic function} on \( I \).
    
    % 或者等价地，实解析函数也可以定义为在定义域内每一点 x_0 的泰勒级数皆逐点收敛的光滑函数 
    Or equivalently, \(f(x)\) is analytic on \(I\) if for any point \(x_{0} \in I\),
    the Taylor series of \(f(x)\) at \(x_{0}\) converges pointwise to \(f(x)\) in some neighborhood of \(x_{0}\), that is,
    \[
    T(x) = \sum_{n=0}^{\infty} \frac{f^{(n)}(x_{0})}{n!} (x - x_{0})^{n} \to  f(x).
    \]

    % I 上实解析函数的全体通常记为
    The set of all real analytic functions on \(I\) is usually denoted by \(C^{\omega}(I)\).
\end{definition}


\section{Smooth Appropriation of Functions}
First, we use continuous functions to approximate Riemann integrable functions
and smooth functions to approximate continuous functions, respectively.
\begin{theorem}
    Let \( f(x) \in R[a, b] \). For any \( \varepsilon > 0 \), 
    there exists a function \( g(x) \in C[a, b] \) such that:
    \[
    \int_{a}^{b}|f(x) - g(x)| < \varepsilon.
    \]
\end{theorem}

\begin{theorem}
    Let \( f(x) \in C[a, b] \). For any \( \varepsilon > 0 \), 
    there exists a function \( g(x) \in C^{\infty}[a, b] \) such that:
    \[
    |f(x) - g(x)| < \varepsilon, \quad \forall x \in [a, b].
    \]
\end{theorem}

Then, Weierstrass approximation theorem is stated as follows:
\begin{theorem}{Weierstrass First Approximation Theorem}
    Let \( f(x) \in C[a, b] \). For any \( \varepsilon > 0 \), 
    there exists a polynomial \( P(x) \) such that:
    \[
    |f(x) - P(x)| < \varepsilon, \quad \forall x \in [a, b].
    \]
\end{theorem}

\begin{theorem}{Weierstrass Second Approximation Theorem}
    Let \(f(x)\) be a continuous periodic function with period \(2\pi\). For any \( \varepsilon > 0 \), 
    there exists a trigonometric polynomial sequence
    \[
    \left\{ T_{n}(x)=\frac{A_{0}}{2} + \sum_{k=1}^{n} A_{k} \cos(kx) + B_{k} \sin(kx) \right\}
    \]
    such that:
    \[
    T_{n}(x) \mathop{\rightrightarrows} f(x).
    \]
\end{theorem}
